\documentclass[10pt,a4paper]{amsart}
\usepackage[margin=19mm]{geometry}
\usepackage{amsmath,amssymb,mathtools}
\usepackage[scr=rsfs,cal=euler]{mathalfa}
\usepackage{xfrac}
\usepackage[unicode,hidelinks,bookmarks=false]{hyperref}
\newcommand{\ie}{i.e.\ }
\newcommand{\cf}{cf.\ }
\newcommand{\resp}{respectively\ }
\newcommand{\Subsection}{\subsection}
\providecommand{\nobreakdash}{\nobreak\hbox{-}}
\setlength{\emergencystretch}{2em}
\multlinegap0pt
\usepackage{amssymb}
\usepackage{bm}
\usepackage{xcolor}
\usepackage[shortlabels]{enumitem}
\usepackage{float}
\newtheorem{theorem}{Theorem}
\DeclareMathOperator*{\Diag}{Diag}
\DeclareMathOperator*{\Id}{Id}
\title{On the monotonicity of the entropy production in the Landau-Maxwell equation}
\author{C\^ome Tabary}
\date{Source: arXiv:2601.03107v5 (CC BY 4.0)}
\begin{document}
\maketitle

\noindent\emph{Source and licence.} On the monotonicity of the entropy production in the Landau-Maxwell equation. Version arXiv:2601.03107v5 (CC BY 4.0), distributed by arXiv under the Creative Commons Attribution 4.0 International licence, http://creativecommons.org/licenses/by/4.0/, as recorded in the arXiv licence metadata for this exact version and shown on the abstract page https://arxiv.org/abs/2601.03107v5. Full licence notice: \texttt{licenses/arxiv-2601.03107-CC-BY-4.0.txt}.\par\smallskip

\noindent\textit{Editorial scope.} Counted unit: the article's theorem thm:main (source0.tex lines 162--178). The article's complete original proof of it is delivered with it (source0.tex lines 649--686, one \texttt{proof} environment of 38 lines, which the article places away from the statement). No mathematics is written, repaired or re-derived: the statement and the proof are the article's own lines, in the article's own notation, and no retained line is substituted or re-flowed.  Retained uncounted as the source-local context the proof needs: lines 115--118; lines 157--160; lines 297--300; lines 607--614.  Nothing was omitted from any retained span, so the package carries no omission record.  No retained line contains a citation, so the wrapper carries no bibliography and no bibliography entry.  No retained line contains a figure environment, an image-inclusion command or drawing code, so no image is delivered and the package declares no asset dependency.  Editorial formatting only, in addition: an AMS \texttt{amsart} preamble that restates the article's own declarations and macros used by the retained lines, namely the environments \texttt{theorem} on separate counters, together with the standard packages those lines need.  The retained environments print in this document as Theorem 1 in order of appearance, whereas the article prints the same items with the numbering of its own full document.  The complete native source of the article as distributed in the arXiv e-print for this exact version is bundled unchanged as \texttt{sources/192207/source0.tex}.
\begin{equation}
\label{eq:landau}
    \partial_t f_t(v) = \nabla_v \cdot \int_{\mathbb{R}^d} \alpha(\vert v-w \vert) a (v - w) (\nabla_v - \nabla_w) f_t(v)f_t(w) dw,
\end{equation}

 \begin{align}
    \label{eq:normalization}
        \int_{\mathbb{R}^d} f_0 dv=1, &&\! \int_{\mathbb{R}^d} v f_0 dv=0,&& \! \int_{\mathbb{R}^d} (v\otimes v) f_0 dv=T_0=\Diag(T_{0,i})_{1\leq i \leq d},&&\!  \sum_{i=1}^d T_{0,i}=d.
    \end{align}

\begin{align}
\label{eq:T}
    T(t) = e^{-4dt/(d-1)} (T_0-\Id)+\Id && \chi(t)=e^{-4dt/(d-1)}\chi_0.
\end{align}

\begin{align}
\label{eq:derivD3}
    \frac{d}{dt}D(f)
    =-&\frac{2(2d+1)}{d-1}I_\chi(f\vert m)+2I_{\Id}(f\vert m)-4I_{(\Id +\chi)^2}(f\vert m)+\mathcal{E}\\
    &-\sum_{i<j,k<l}K^{ \mathcal{B}_{ij}\mathcal{B}_{kl}}(f\vert m)-2\sum_i\sum_{k<l}(1+\chi_i)K^{\partial_i\mathcal{B}_{kl} }(f\vert m)\nonumber\\
    &-\sum_{i,j}(1+\chi_i)(1+\chi_j)K^{\partial_i\partial_j}(f\vert m)\nonumber\\
    &-(12d-4)\Vert \chi\Vert^2-4(d-1)\sum_i\chi_i^3,\nonumber
\end{align}

\begin{theorem}
\label{thm:main}
    Let $f_0\geq 0$ with finite mass and energy satisfy the normalization \eqref{eq:normalization}
    and let $T_{0,\max}=\max_{1\leq i \leq d} (T_{0,i})\geq 1$. Consider $f=(f_t)_{t\geq 0}$ the solution to the Landau equation \eqref{eq:landau} with Maxwell molecules $\alpha=1$ and initial condition $f_0$.
    Then, there exists $t_0 \geq 0$ depending only on $d$ and $T_{0,\max}$ such that for all $t\geq t_0$, $D(f_t)$ is monotone non-increasing.
    We can further take 
    $$t_0=\frac{d-1}{4d}\log\left(\frac{T_{0,\max}-1}{C_d}\right)$$
    or $t_0=0$ if the above quantity is negative. The constant $C_d$ is given by
    \begin{align*}
       C_d= \begin{cases}
            \frac{13-\sqrt{133}}{18}\approx 0.08 & \text{if }d=2 \\
            \frac{20-2\sqrt{70}}{15}\approx 0.2 & \text{if }d=3 \\
         (d-1)^2\left((6d-3)+\frac{d(d+1)}{\eta(d-1)}\right)^{-1}\approx_{d\rightarrow \infty} 0.15d & \text{for }d\geq 4,
        \end{cases}
    \end{align*}
    where $\eta=\frac{9+\sqrt{129}}{12}\approx 1.69$.
\end{theorem}

\begin{proof}[Proof of Theorem~\ref{thm:main}]
We start from the relative form of $\frac{d}{dt}D(f)$ expressed in \eqref{eq:derivD3}. We drop two out of the three non-negative $K$ terms, and plug in the result of the previous lemma for some $\eta>0$ to fix later:
\begin{align}
\label{eq:derivD4}
    \frac{d}{dt}D(f)
    \leq-&\frac{2(2d+1)}{d-1}I_\chi(f\vert m)+2\left[1+\frac{\eta^{-1}}{d-1}\Vert \chi\Vert^2\right]I_{\Id}(f\vert m)-4I_{(\Id +\chi)^2}(f\vert m) \\
    &-2\sum_i\sum_{k<l}(1+\chi_i-\eta/2)K^{\partial_i\mathcal{B}_{kl} }(f\vert m)\nonumber\\
    &+\left[\frac{2d\eta^{-1}}{d-1}-(12d-4)\right]\Vert \chi\Vert^2-4(d-1)\Vert \chi\Vert^3\nonumber.
\end{align}
Since $\vert \chi_i\vert\leq 1$, we have $\vert \sum_i \chi_i^3 \vert \leq \Vert \chi\Vert^2$, and the last line is bounded by
$$\left[\frac{2d\eta^{-1}}{d-1}-(12d-4)\right]\Vert \chi\Vert^2-4(d-1)\Vert \chi\Vert^3 \leq \left[\frac{2d\eta^{-1}}{d-1}-8d\right]\Vert \chi\Vert^2.$$
If we pick $\eta\geq \frac{1}{4(d-1)}$, the right-hand side is non-positive, and the last line of \eqref{eq:derivD4} is non-positive for all times. For the second line, we need to pick $\eta<2$, and it is then non-positive as soon as
\begin{align*}
    \max_i(T_i)-1=(d-1)\max_i(-\chi_i)\leq (d-1)(1-\eta /2)=:C_1(d,\eta)
\end{align*}
The first line is non-positive as soon as the matrix
\begin{align*}
    -\left[\frac{2(2d+1)}{d-1}+
    8\right]\chi +\left[-2+\frac{2\eta^{-1}}{d-1}\Vert \chi\Vert^2\right] \Id -4\chi^2\leq 0.
\end{align*}
This always happens since the left-hand side goes to $-2\Id$ as $\chi\rightarrow 0$. To get a quantitative bound, since $$\Vert\chi\Vert^2=\frac{1}{(d-1)^2}\left(\sum_i T_i^2-d\right)\leq \frac{d}{(d-1)^2}\left((\max_i(T_i))^2-1\right)\leq \frac{d(d+1)}{(d-1)^2}\left((\max_i(T_i))-1\right),$$ we can crudely bound every (diagonal) entry of the left-hand side by:
$$ \left[\frac{2(2d+1)}{d-1}+
    8\right]\frac{\max_i(T_i)-1}{d-1} +\frac{2\eta^{-1}d(d+1)}{(d-1)^3}\left( (\max_i(T_i))-1\right)-2.$$
Multiplying by $(d-1)^2$, we see that this is non-positive as soon as
$$\max_i(T_i)-1\leq (d-1)^2\left((6d-3)+\frac{\eta^{-1}d(d+1)}{d-1}\right)^{-1} =:C_2(d,\eta).$$
We now optimize $\eta\in[\frac{1}{4(d-1)},2)$ to get matching constants $C_1$ and $C_2$.
\begin{itemize}
    \item For $d=2$, we can pick $\eta=\frac{5+\sqrt{133}}{9}\approx 1.83$ so that $C_1(2,\eta)=C_2(2,\eta)=\frac{13-\sqrt{133}}{18}$. 
    \item For $d=3$, we can pick
$\eta=\frac{10+2\sqrt{70}}{15}\approx 1.78$ so that $C_1(3,\eta)=C_2(3,\eta)=\frac{20-2\sqrt{70}}{15}$.
    \item As $d\rightarrow \infty$, $C_1(d,\eta) \sim d(1-\eta/2) $ and $C_2(d,\eta) \sim d/(6+\eta^{-1}) $, so we take $\eta=\frac{9+\sqrt{129}}{12}\approx 1.69$, for which $C_1(d,\eta)\sim C_2(d,\eta)$, and for which one can check that $C_1(d,\eta)\geq C_2(d,\eta)$ always.
\end{itemize}
In any case, the entropy production is non-increasing as soon as $\max_i(T_i)-1\leq C_2(d,\eta)$. If this is not the case initially, the evolution of $T$ given by \eqref{eq:T} tells us that it is guaranteed as soon as
\begin{align*}
    e^{-4dt/(d-1)}(T_{0,\max}-1)\leq C_2(d,\eta), && \textit{i.e.} && t\geq \frac{d-1}{4d}\log\left(\frac{T_{0,\max}-1}{C_2(d,\eta)}\right).
\end{align*}
This proves Theorem~\ref{thm:main} with $C_d=C_2(d,\eta)$.
\end{proof}

\end{document}
